This thesis investigates automatic, style- and data-driven co-optimisations for MapLibre-based map rendering pipelines.
Interactive web maps built on vector tiles face a complex performance landscape: every frame requires fetching tiles over the network, decoding them, evaluating style expressions for each feature, and submitting draw calls to the graphics API.
Despite this, style authoring and tile data preparation have traditionally been treated as independent concerns, leaving the interaction between the two unexploited.

This work designs and implements a standalone optimiser that operates at three levels.
Expression-level passes apply a set of rewrite rules - constant folding, algebraic simplification, stats-driven folding from tile statistics, and selectivity-based predicate reordering - to the style's JSON expressions via fixpoint iteration.
Structural passes then transform typed representations of the style document, performing dead layer elimination, metadata refinement, layer merging, and cleanup.
Beyond the style, the optimiser bridges into data-level optimisations through style-driven tile shaving, which computes a pruning advisory from the optimised style to remove unused properties, geometry types, and feature values from tile data before re-encoding into the columnar MapLibre Tiles format with automatic encoding strategy selection.

An evaluation across 15~styles, 18~scenarios, and 17~cumulative ablation steps demonstrates that the full pipeline achieves load-time improvements of \qtyrange{71}{75}{\percent} for the two styles benchmarked end-to-end (Fiord and Liberty; the remaining styles are evaluated at the style level only) and boosts FPS by \qtyrange{162}{213}{\percent} through tile shaving.
The Rust-based \ac{MLT} decoder achieves \numrange{4.6}{8.3}$\times$ higher decode throughput than \ac{MVT}+gzip, while the encoder reduces total tile volume by \qty{28}{\percent} (\qty{4.24}{\giga\byte} to \qty{3.06}{\giga\byte} on the Germany OpenMapTiles tileset), \qty{28}{\percent} over the reference Java encoder, and \qty{12}{\percent} over gzip-compressed \ac{MVT}.
Style-level optimisations alone reduce style transfer size by \qtyrange{3}{47}{\percent} (gzip), primarily benefiting initial load.
Style-level passes simplify the advisory computation and provide modest rendering improvements; tile shaving and encoding strategy selection deliver the dominant load-time and transfer-size gains.
The two levels compose additively rather than synergistically: the aggregate data shows no super-additive interaction at the median level, though per-scenario synergy is present in a subset of configurations.
